\subsection{Rozwiązanie}
\paragraph{Zadanie programowania liniowego}

\begin{lstlisting}[caption= plik dat]
data; 
 
set BudynkiNaStart := s, A, B, C, D, E, F, G, H; 
set Budynki := A, B, C, D, E, F, G, H; 
set Elektrownie := F, G, H; 

param Z := F 15 G 13 H 7; 
 
param p := 
s A 10 s B 13 s C 17 s D 0  s E 0  s F 0  s G 0  s H 0
A A 0  A B 0  A C 0 A D 15  A E 10 A F 0  A G 0  A H 0
B A 0  B B 0  B C 0 B D 4 B E 9  B F 0  B G 9  B H 0
C A 0  C B 0 C C 0 C D 20 C E 10 C F 0  C G 0  C H 0   
D A 0  D B 0  D C 0  D D 0 D E 0 D F 10 D G 3   D H 2 
E A 0  E B 0  E C 0 E D 20 E E 0 E F 5 E G 5  E H 5  
F A 0  F B 0  F C 0  F D 0  F E 0  F F 0  F G 0  F H 0
G A 0  G B 0  G C 0 G D 0  G E 0  G F 0  G G 0  G H 0
H A 0  H B 0  H C 0  H D 0  H E 0  H F 0  H G 0 H H 0; 
 
param c := 
s A 0 s B 0 s C 0 s D 0 s E 0 s F 0 s G 0 s H 0
A A 0 A B 0 A C 0 A D 4  A E 2 A F 0 A G 0 A H 0
B A 0 B B 0 B C 0 B D 4 B E 3 B F 0 B G 8 B H 0
C A 0 C B 0 C C 0 C D 2 C E 6 C F 0 C G 0 C H 0   
D A 0 D B 0 D C 0 D D 0 D E 0 D F 3 D G 7 D H 2
E A 0 E B 0 E C 0 E D 5  E E 0 E F 7 E G 6 E H 3 
F A 0 F B 0 F C 0 F D 0 F E 0 F F 0 F G 0 F H 0
G A 0 G B 0 G C 0 G D 0 G E 0 G F 0 G G 0 G H 0
H A 0 H B 0 H C 0 H D 0 H E 0 H F 0 H G 0 H H 0; 

end; 
\end{lstlisting}
\begin{lstlisting}[caption= plik mod]
set BudynkiNaStart; 
set Budynki; 
set Elektrownie; 

param Z{e in Elektrownie}; 
param p{i in BudynkiNaStart, j in Budynki}; 
param c{i in BudynkiNaStart, j in Budynki}; 

var f{i in BudynkiNaStart, j in Budynki}, >= 0; 
 
minimize Q: sum {i in Budynki, j in Budynki} c[i,j] * f[i,j]; 
 
subject to 
  Ograniczenie_jeden{i in BudynkiNaStart, j in Budynki}: 
    f[i,j] >= 0; 
  Ograniczenie_dwa{i in BudynkiNaStart, j in Budynki}: 
    f[i,j] <= p[i,j]; 
  Ograniczenie_trzy{e in Elektrownie}: 
      sum {n in BudynkiNaStart} f[n,e] >= Z[e]; 
  Ograniczenie_cztery{n in Budynki}: 
    sum {i in Budynki} f[n,i] <= sum {j in BudynkiNaStart} f[j,n]; 
 
solve; 
display {i in Budynki, j in Budynki: f[i,j] > 0}: f[i,j]; 
display: Q;
\end{lstlisting}
\paragraph{Wynik}
\begin{lstlisting}[caption=Wynik z glpk]
f[A,E].val = 10
f[B,E].val = 4
f[B,G].val = 9
f[C,D].val = 12
f[D,F].val = 10
f[D,H].val = 2
f[E,F].val = 5
f[E,G].val = 4
f[E,H].val = 5
Q.val = 236
\end{lstlisting}

\begin{figure}[htb]
\caption{Plan dostaw węgla}
\begin{center}
\begin{tabular}{ | c | c | c | c | c | c | }
\hline
  & D & E & F & G & H \\ 
 \hline
 A & 10 & - & - & - & -\\  
 \hline
 B & 0 & 4 & - & 9 & - \\
 \hline  
 C & 12 & - & - & - & - \\
 \hline  
 D & - & - & 10 & - & 2 \\
 \hline  
 E & - & - & 5 & 4 & 5 \\
 \hline  
\end{tabular}
\end{center}
\end{figure}

\[ Q = 236 \]

